\documentclass[thesis=master, language=german, paper=a4]{isys}[2013/06/28]
%
% [2013/06/28] ensure that the newest isys.cls is used
%              if an older version is used, latex issues a warning
%
% thesis   = sa       - Studienarbeit  / Student thesis 
%            da       - Diplomarbeit   / Diploma thesis
%            bachelor - Bachelorarbeit / Bachelor thesis
%            master   - Masterarbeit   / Master thesis
%            semester - Semesterarbeit / Semester thesis
%
% language = {english, german} 
%            Set default language of document. The class file isys.cls loads the babel
%              package with both options [english, ngerman]. For text that shall be 
%              typeset in a different language than the default language, the 
%              environment 'otherlanguage' or the command '\foreignlanguage' may be used
%              with either of these options (see babel package documentation)
%
% paper    = {a4, letter} - paper size
%
% thesisdraft - mark document as draft without passing "draft" option to 
%                 other packages (e.g. pictures are shown in the document, etc.)
%
% oneside - this option avoids empty pages in the document output 
%                which are necessary for twoside printing, which is useful for 
%                draft versions. 
%                DO NOT USE THIS OPTION WHEN SUBMITTING YOUR THESIS!
%

% mark overful boxes by a 2pt wide rule 
% -- FIX ALL OVERFULL BOXES BEFORE SUBMITTING YOUR THESIS!
\overfullrule=2pt 


%%%%%%%%%%%%%%%%%%%%%%%%%%%%%%%%%%%%%%%%%%%%%%%%%%%%%%%%%%%%%%%%%%%
% Load packages - useful packages, adapt to your own needs
%
% isys.cls loads the "graphicx" package, so \includegraphics can 
%  directly be used to include figures.
%
%%%%%%%%%%%%%%%%%%%%%%%%%%%%%%%%%%%%%%%%%%%%%%%%%%%%%%%%%%%%%%%%%%%
% input encoding - UTF8 or latin1 allow directly typing german umlauts
% make sure the file encoding of all files matches this encoding!
\usepackage[utf8]{inputenc}
% font enconding - do not change this unless you exactly know what you're doing!
\usepackage[T1]{fontenc}

% automatically use correct quotes for the current language: \enquote{text in quotes}
\usepackage{csquotes}
\MakeAutoQuote{»}{«}  % input »text in quotes« is equivalent to \enquote{text in quotes}
\let\q=\enquote 			% use \q{text in quotes}, which is shorter than \enquote{text in quotes}

% better math-commands (e.q.: \dddots, \eqref{})
\usepackage{amsmath,amssymb}                                 
\usepackage{bm}
% better-looking rules in tabulars, \toprule, \midrule, \bottomrule
\usepackage{booktabs}

% use consistent style for units (upright shape)
% \unit[value]{unit}, \unitfrac[value]{unit_numerator}{unit_denominator}
\usepackage{units}

% tables with fixed table width 
\usepackage{tabularx, tabulary}
% improve ragged right alignment with \RaggedRight
\usepackage{ragged2e}
% enable import of eps pictures using pdflatex
%\usepackage{epstopdf}

% enable hyperlinks and bookmarks in PDF file (must typically be loaded as LAST package)
% if used, print hyperlinks in black and suppress boxes around them 
\usepackage[colorlinks=true, allcolors=black, bookmarksnumbered=true]{hyperref}
% to set bookmarks on the first latex run (hyperref needs two),
% provides addtional options for bookmarks (use carefully!)
\usepackage{bookmark}

% für mehrere Grafiken nebeneinander
\usepackage{subcaption}
\usepackage{svg} % um SVGs direkt einbinden zu können
\usepackage{tikz}
\usepackage{pgfplots} 
\usepackage{pgfgantt}
\usepackage{pdflscape}
\pgfplotsset{compat=newest} 
\pgfplotsset{plot coordinates/math parser=false}

%% the following commands are needed for some matlab2tikz features
\usetikzlibrary{plotmarks}
\usetikzlibrary{arrows.meta}
\usepgfplotslibrary{patchplots}
\usepackage{grffile}

%% um eps Dateien einzubinden
\usepackage{graphicx}
\usepackage{import} 

% Bibliography using classical bibtex
%\bibliographystyle{alphadin}

% better: use biblatex! provides much more options
% use \printbibliography to print the bibliography 
% (search for \printbibliography in this document and uncomment the line, 
%  comment out the line for bibtex)
\usepackage[style=numeric,
						backend=bibtex8,		% or use biber or bibtex?
						bibencoding=utf8,	% .bib file in latin1 or utf8 encoding?
						maxbibnames=99			% give full list of authors (up to 99) in bibliography
						]{biblatex}
\bibliography{references}


%%%%%%%%%%%%%%%%%%%%%%%%%%%%%%%%%%%%%%%%%%%%%%%%%%%%%%%%%%%%%%%%%%%
% Information for title page 
%%%%%%%%%%%%%%%%%%%%%%%%%%%%%%%%%%%%%%%%%%%%%%%%%%%%%%%%%%%%%%%%%%%

% title in the language in which the thesis is written
\title{Schätzverfahren für adaptive Regelung von Hybridgetrieben}
%\subtitle{ }
% title in foreign language (german or english, depending on the language of the thesis)
\foreigntitle{Estimation Methods for Adaptive Control of Hybrid Gearboxes}
%\foreignsubtitle{Estimation Procedure for Adaptive Control of Hybrid Gearboxes}

\author{Koray Karci}
\major{Maschinenbau}
\date{30.09.2018}		% use either ISO format YYYY-MM-DD or german format DD.MM.YYYY
										% default: current date

\examiner{Prof. Dr.-Ing. Oliver Sawodny}
%\examiner{Prof. Dr.-Ing. Cristina Tar\'{i}n}
%\supervisor{Dipl.-Ing. Supervisor 1}		% give max. 4 \supervisor {...} commands
\supervisor{M.Sc. Michel Bauer}
%\supervisor{Dipl.-Ing. Supervisor 3}
%\supervisor{M.Sc. Supervisor 4}

%% if work on thesis has been conducted with a partner company/institute
%% use \partner and (if desired also \partnerLogo)
% \partnerLogo has no effect as long as no \partner{...} is given
%\partner{Example company name\\
%			Second line company name}
%\partnerLogo{\includegraphics[height=1.4cm, clip, viewport=0 0 30 10]{logos/logoISYS_deutsch}}



%%%%%%%%%%%%%%%%%%%%%%%%%%%%%%%%%%%%%%%%%%%%%%%%%%%%%%%%%%%%%%%%%%%
% Start of document output
%%%%%%%%%%%%%%%%%%%%%%%%%%%%%%%%%%%%%%%%%%%%%%%%%%%%%%%%%%%%%%%%%%%
\begin{document}

% Titelblatt ausgeben / print title page
\pdfbookmark{Titelseite}{title}
\maketitle

%% Englisches Titelblatt in deutschsprachiger Arbeit erstellen
%\begin{otherlanguage}{english}\maketitle \end{otherlanguage}

%% Create german title page in english theses
%\begin{otherlanguage}{ngerman}\maketitle \end{otherlanguage}


	
% Kurzfassung in der Dokumentensprache / Abstract in default document language
\pdfbookmark{Kurzfassung}{abstract-de}
\begin{abstract}
Kurzfassung
\end{abstract}

% Falls gewuenscht, englischen Abstract in deutschsprachiger Arbeit einbinden
\newpage
\begin{otherlanguage}{english}
\pdfbookmark{Abstract}{abstract-en}
\begin{abstract}
Abstract 
\end{abstract}
\end{otherlanguage}

%% Inhalts-/Abbildungs-/Tabellenverzeichnis
%% Table of contents, List of Figures/Tables

	
	
%% If subsections shall also be included in the table of contents
%% uncomment next line. Default: only sections are included.
\setcounter{tocdepth}{1} % include only chapters and sections (1.1) - default
%\setcounter{tocdepth}{2} % also include subsections - use only if told to do so by your advisor!
\cleardoublepage
\pdfbookmark{\contentsname}{toc}
\tableofcontents

\addchap{Abkürzungen und Formelzeichen}

\section*{Abkürzungen}

\begin{flushleft}
\begin{tabularx}{\textwidth}{XlX}	
	&&\\
	DOF								&Freiheitsgrade (Engl.: Degrees of Freedom)&\\
	E/A								&Eingang/Ausgang &\\
	IMC								&Internal Model Control&\\
	MPC								&Modellprädiktive Regelung (Engl.: Model Predictive Control)&
\end{tabularx}
\end{flushleft}


\vspace{20mm}
\section*{Formelzeichen}


\begin{flushleft}
\begin{tabularx}{\textwidth}{llX}
	\toprule
	Formelzeichen					 	& Einheit									& Beschreibung \\
	\midrule
	$I$											&	A												&	Strom \\
	$U$											&	V												&	Spannung \\
	$T$											&	K												&	Temperatur \\
	&&\\
	$\tau$									&	s												& Zeitkonstante \\
	$\rho_{\mathrm{L}}$			&	$\unitfrac[]{kg}{m^3}$	& Dichte der Luft \\
	&&\\
	$\boldsymbol{u}$						& 												&	Regeleingang des zu regelnden Systems \\
	$\boldsymbol{x}$						& 												&	Zustandsvektor \\
	$\boldsymbol{y}$						& 												&	Regelausgang des zu regelnden Systems \\
\bottomrule
\end{tabularx}
\end{flushleft}


%% Chapters and sections without numbering but included in running heading
%\addchap{Chapter without numbering, included in table of contents}
%\addchap*{Chapter without numbering, not included in table of contents}
%\addsec{Section without numbering, included in table of contents}
%\addsec*{Section without numbering, not included in table of contents}

%%%%%%%%%%%%%%%%%%%%%%%%%%%%%%%%%%%%%%%%%%%%%%%%%%%%%%%%%%%%%%%%%%%
% Beginn des Textes Inhalts / Start of thesis content
%%%%%%%%%%%%%%%%%%%%%%%%%%%%%%%%%%%%%%%%%%%%%%%%%%%%%%%%%%%%%%%%%%%

	
% It may be helpful to put each chapter into a separate file. 
% \includeonly{...} allows to restrict the files that are included,
% which may speed up the document creation process.
% \include always starts a new odd page before including the file 
% contents. If this is not desired, \input can be used.
\chapter{Einleitung und Motivation}\label{ch:ch1}
Dies ist die Einleitung und Motivation.




\chapter{Modellierung}\label{ch:ch2}
In diesem Kapitel wird ein für die Simulation verwendetes detailliertes Modell des Antriebsstrangs hergeleitet. Hierfür wird ein vorhandenes, starres Modell um Elastizitäten der Seitenwellen ergänzt. Für die Herleitung wird dabei der Lagrange Formalismus verwendet, welcher daher auch kurz beschrieben wird. Im weiteren werden aus dem detaillierten Modell vereinfachte Modelle hergeleitet, welche dem Entwurf der Beobachter im Kapitel~\ref{ch:ch4} dienen.


\section{Modellierung des Antriebsstrangs}
Zur Modellierung des Antriebsstranges wird zunächst die Getriebedynamik mittels der Newton-Euler Gleichungen berechnet. Dazu wird zu Beginn kurz auf die Herleitung dieser Gleichungen eingegangen. Danach wird dieses um die Elastizitäten der Seitenwellen ergänzt, woraus sich ein schwingfähiges System ergibt.  

\subsection{Beschreibung des allgemeinen Newton-Euler-Formalismus}
Die Position und Orientierung von jedem Körper $i$ eines holonomen Mehrkörpersystems mit $p$ Körpern und $f$ Freiheitsgraden (FHG), kann in Abhängigkeit der verallgemeinerten Koordinaten $\pmb{y}\in \mathbb{R}^f$ angegeben werden mit
\begin{align}
\pmb{r}_i &= \pmb{r}_i(\pmb{y},t)\quad i = 1(1)p\\
\pmb{s}_i &= \pmb{s}_i(\pmb{y},t).
\end{align}
Durch die zeitliche Ableitung dieser ergeben sich die translatorischen und rotatorischen Geschwindigkeiten
\begin{align}\label{vi}
\pmb{v}_i &=  \dot{\pmb{r}}_i = \frac{\partial\pmb{r}_i}{\partial\pmb{y}}\cdot \dot{\pmb{y}} + \frac{\partial\pmb{r}_i}{\partial\pmb{t}} = \pmb{J}_{Ti}(\pmb{y},t)\cdot \dot{\pmb{y}} + \overline{\pmb{v}}_i(\pmb{y},t)\\\label{omegai}
\pmb{\omega}_i &=  \dot{\pmb{s}}_i = \frac{\partial\pmb{s}_i}{\partial\pmb{y}}\cdot \dot{\pmb{y}} + \frac{\partial\pmb{s}_i}{\partial\pmb{t}} = \pmb{J}_{Ri}(\pmb{y},t)\cdot \dot{\pmb{y}} + \overline{\pmb{\omega}}_i(\pmb{y},t)
\end{align}
mit der Jacobi-Matrix der Translation $\pmb{J}_{Ti}$ und der Rotation $\pmb{J}_{Ri}$. Die lokalen Geschwindigkeitsvektoren $\overline{\pmb{v}}_i(\pmb{y},t)$ und $\overline{\pmb{\omega}}_i(\pmb{y},t)$ treten nur bei Systemen mit rheonomen Bindungen auf. Nach erneuter zeitlicher Ableitung ergeben sich die Beschleunigungen 
\begin{align}
\pmb{a}_i = \dot{\pmb{v}}_i &= \pmb{J}_{Ti} \cdot \ddot{\pmb{y}} + \dot{\pmb{J}}_{Ti} \cdot \dot{\pmb{y}} + \dot{\overline{\pmb{v}}}_i = \pmb{J}_{Ti}(\pmb{y},t) \cdot \ddot{\pmb{y}} + \overline{\pmb{a}}_i(\pmb{y},\dot{\pmb{y}},t)\\
\pmb{\alpha}_i = \dot{\pmb{\omega}}_i &= \pmb{J}_{Ri} \cdot \ddot{\pmb{y}} + \dot{\pmb{J}}_{Ri} \cdot \dot{\pmb{y}} + \dot{\overline{\pmb{\omega}}}_i = \pmb{J}_{Ri}(\pmb{y},t) \cdot \ddot{\pmb{y}} + \overline{\pmb{\alpha}}_i(\pmb{y},\dot{\pmb{y}},t)\label{alphai}
\end{align}
mit den lokalen Beschleunigungsvektoren $\overline{\pmb{a}}_i$ und $ \overline{\pmb{\alpha}}_i$. 
Des weiteren lassen sich für die Dynamik eines Systems virtuelle Bewegungen definieren. Diese sind willkürliche, infinitesimale Bewegungen des Systems, welche mit den skleronomen und den rheonomen Bindungen verträglich sind. Für virtuelle Bewegungen an holonomen Bindungen gilt
\begin{align}
\delta \pmb{r} &\neq \pmb{0},\\
\delta \pmb{r} &= \pmb{0},\label{VB0}\\
\delta t &= 0.
\end{align} 
Damit lassen sich die virtuellen Bewegungen der einzelnen Körper $i$ definieren zu
\begin{align}\label{VBr}
\delta \pmb{r}_i = \frac{\partial\pmb{r}_i}{\partial\pmb{y}}\cdot \delta\pmb{y} &= \pmb{J}_{Ti}(\pmb{y},t)\cdot\delta\pmb{y}\\
\delta \pmb{s}_i = \frac{\partial\pmb{s}_i}{\partial\pmb{y}}\cdot \delta\pmb{y} &= \pmb{J}_{Ri}(\pmb{y},t)\cdot\delta\pmb{y}.\label{VBs}
\end{align}
Die Kinetik eines Körpers $i$ wird durch die Newtonschen- und Eulerschen-Gleichungen beschrieben. 
Die Newtonschen-Gleichungen eines Körpers $i$ lassen sich angeben als
\begin{equation}\label{NewtonGl}
m_i \pmb{a}_i(t) = \pmb{f}_i(t)
\end{equation}
wobei $m_i$ die Masse des Körpers angibt und $\pmb{f}_i(t)$ die Summe der angreifenden Kräfte. Diese lassen sich wiederum einteilen in eingeprägte Kräfte $\pmb{f}^e$ und Reaktionskräfte $\pmb{f}^r$.
Die Eulerschen Gleichungen lauten
\begin{equation}\label{EulerGl}
\pmb{I}_i(t)\cdot\pmb{\alpha}_i(t) + \tilde{\pmb{\omega}}_i(t)\cdot \pmb{I}_i(t)\pmb{\omega}_i(t)= \pmb{l}_i(t)
\end{equation}
mit dem Trägheitstensor $\pmb{I}_i(t)$ und den äußeren Momenten $\pmb{l}_i(t)$ bezogen auf den Massenmittelpunkt des Körpers $i$ im Inertialsystem und dem schiefsymmetrische 3x3-Tensor $\tilde{\pmb{\omega}}_i(t)$. Die äußeren Momenten $\pmb{l}_i(t)$ lassen sich äquivalent zum translatorischen Fall wieder in eingeprägte Momente $\pmb{l}^e$ und Reaktionsmomente $\pmb{l}^r$ aufteilen. Der schiefsymmetrische 3x3-Tensor ist für ein beliebigen Vektor 
\begin{equation}
\pmb{v} = [v_1 \quad v_2 \quad v_3]^T
\end{equation}
definiert als
\begin{equation}
\tilde{\pmb{v}} = \begin{bmatrix} 0 & -v_3 & v_2 \\ v_3 & 0 & -v_1 \\ -v_2 & v_1 & 0 \end{bmatrix}.
\end{equation}
Darüber hinaus kann mit den virtuellen Bewegungen aus (\ref{VBr}) und (\ref{VBs}) die virtuelle Arbeit $\delta \pmb{W}$ definiert werden. Da es, wie in (\ref{VB0}) definiert, in Richtung der Bindungen keine virtuellen Bewegungen geben kann, ergibt sich für die virtuelle Arbeit der Reaktionskräfte
\begin{equation}\label{VA}
\delta\pmb{W}^r = \sum_{i=1}^p(\pmb{f}_i^r\cdot \delta \pmb{r}_i + \pmb{l}^r_i\cdot\delta\pmb{s}_i)=0.
\end{equation}
Mit den Gleichungen \ref{NewtonGl}, \ref{EulerGl} und der beschriebenen Aufteilung der äußeren Kräfte in eingeprägte Kräfte und Reaktionskräfte, folgt aus (\ref{VA}) das d’Alembertsche Prinzip für Mehrkörpersysteme
\begin{equation}\label{VA}
\sum_{i=1}^p\left[ (m_i\pmb{a}_i-\pmb{f}_i^e)\cdot \delta \pmb{r}_i + (\pmb{I_i}\cdot \pmb{\alpha}_i + \tilde{\pmb{\omega}}\cdot\pmb{I}_i\cdot\pmb{\omega}_i-\pmb{l}^e_i)\cdot\delta\pmb{s}_i\right] =0.
\end{equation}
Durch einsetzten der virtuellen Bewegungen (\ref{VBr}) und (\ref{VBs}) ergibt sich daraus
\begin{equation}
\delta\pmb{y}\sum_{i=1}^p\left[\pmb{J}^T_{Ti}\cdot(m_i\pmb{a}_i-\pmb{f}_i^e)+ \pmb{J}^T_{Ri}\cdot(\pmb{I_i}\cdot \pmb{\alpha}_i + \tilde{\pmb{\omega}}\cdot\pmb{I}_i\cdot\pmb{\omega}_i-\pmb{l}^e_i)\right]=0.
\end{equation}
Das einarbeiten der kinematischen Gleichungen (\ref{vi}) -- (\ref{alphai}) in ... führt zur Bewegungsgleichung für holonome Mehrkörpersysteme (MKS) 
\begin{align}\label{Bwg_lang}
\begin{split}
\sum_{i=1}^p\left[\pmb{J}_{Ti}^T\cdot m_i \cdot \pmb{J}_{Ti} + \pmb{J}_{Ri}^T\cdot \pmb{I}_i \cdot \pmb{J}_{Ri}\right]\ddot{\pmb{y}}+\sum_{i=1}^p\left[\pmb{J}_{Ti}^T\cdot m_i \cdot \overline{\pmb{a}}_i + \pmb{J}^T_{Ri}\cdot\pmb{I}_i\cdot\overline{\pmb{\alpha}}_i+\pmb{J}^T_{Ri}\cdot\tilde{\pmb{\omega}}_i\cdot\pmb{I}_i\cdot\pmb{\omega_i}\right]=\\ \sum_{i=1}^p\left[\pmb{J}^T_{Ti}\cdot\pmb{f}_i^e + \pmb{J}_{Ri}^T\cdot\pmb{I}_i^e\right].
\end{split}
\end{align}
Darin lässt sich die erste Summe zur Massenmatrix $\pmb{M}(\pmb{y},t)$, die zweite Summe zum Vektor der verallgemeinerten Zentrifugal-- und Coriloriskräfte $\pmb{k}(\pmb{y},\dot{\pmb{y}},t)$ und die dritte Summe zum Vektor der verallgemeinerten Kräfte $\pmb{q}(\pmb{y},\dot{\pmb{y}},t)$ zusammenfassen, wodurch sich (\ref{Bwg_lang}) schreiben lässt als
\begin{equation}\label{Bwg}
\pmb{M}(\pmb{y},t)\cdot\ddot{\pmb{y}}+\pmb{k}(\pmb{y},\dot{\pmb{y}},t) = \pmb{q}(\pmb{y},\dot{\pmb{y}},t)
\end{equation} 
  
\subsection{Beschreibung des Getriebeaufbaus}
Mit der allgemeinen Bewegungsgleichung für MKS (\ref{Bwg}) kann die Getriebedynamik berechnet werde. Das in dieser Arbeit betrachtete Getriebe besteht im wesentlichen aus vier hintereinander liegenden Umlaufgetrieben. Der schematische Aufbau eines Umlaufgetriebes ist in Abbildung ??? zu sehen. Dieses besteht aus dem Sonnenrad in der Mitte, den darum gleichmäßig verteilten Planetenrädern, welche auf dem Planetenträger drehbar gelagert sind und dem Hohlrad. Der Zusammenhang der einzelnen Drehzahlen ist gegeben durch die beiden Zwangsbedingungen (ZB)
\begin{align}
\mathrm{Sonne - Planeten}:\quad &r_{\mathrm{S},j}\omega_{\mathrm{S},j} - r_{\mathrm{T},j}\omega_{\mathrm{T},j}+ r_{\mathrm{P},j}\omega_{\mathrm{P},j}=0 \\
\mathrm{Planeten - Hohlrad}:\quad &r_{\mathrm{H},j}\omega_{\mathrm{H},j} - r_{\mathrm{T},j}\omega_{\mathrm{T},j}- r_{\mathrm{P},j}\omega_{\mathrm{P},j}=0
\end{align}
mit den Rollradien $r_{S,j}$, $r_{P,j}$, $r_{T,j}$ und $r_{H,j}$ des Sonnenrads, des Planetenrads, des Planententrägers und des Hohlrades im $j$-ten Planetensatz. Entsprechenden gelten die Indizes auch für die Winkelgeschwindigkeiten $\omega_{S,j}$, $\omega_{P,j}$, $\omega_{T,j}$ und $\omega_{H,j}$
Dabei gilt, dass die Verhältnisse der Rollradien denen der Zähnezahlen einer Zahnradverbindung entsprechen. Der Rollradius des Planetenträger entspricht somit der Summe der Zähnezahl der Sonne $z_{S,j}$ und des Planeten $z_{P,j}$ und es ergibt sich
\begin{align}\label{ZB_zahnS}
\mathrm{Sonne - Planeten}:\quad &z_{\mathrm{S},j}\omega_\mathrm{S} - (z_{\mathrm{S},j} + z_{\mathrm{P},j})\omega_{\mathrm{T},j}+ z_{\mathrm{P},j}\omega_{\mathrm{P},j}=0 \\ \label{ZB_zahnH}
\mathrm{Planeten - Hohlrad}:\quad &z_{\mathrm{H},j}\omega_{\mathrm{H},j} - (z_\mathrm{S} + z_{\mathrm{P},j})\omega_{\mathrm{T},j}- z_{P,j}\omega_{\mathrm{P},j}=0.
\end{align}
Die Größe  $z_{S,j}$ entspricht der Zähnezahl am Hohlrad im $j$-ten Planetensatz. Die Verschaltung der vier Planetensätze zeigt Abbildung ???. Aus den acht Wellen ergeben sich entsprechende acht FHG im des MKS. Vier weitere kommen durch die Planten der einzelnen Planetensätze hinzu. Somit hat des freie MKS $f_{\textrm{free}}=12$ FHG. Unter Berücksichtigung der beiden Zwangsbedinungen (\ref{ZB_zahnS}) und (\ref{ZB_zahnH}) ergibt sich somit die Anzahl der FHG des beschränkten MKS zu
\begin{equation}
f=f_{\textrm{free}}-4q = 12-8 = 4
\end{equation}
mit der Anzahl der ZB in jedem Planetensatz $q$. Weiterhin ergeben sich durch die ZB an den einzelnen Planetensätzen die Zusammenhänge
\begin{subequations}\label{eq:zwangsbedingungen}
	\begin{align} 
\mathrm{Planetensatz\ 1}: \quad z_{\mathrm{S},1} \omega_{1}  &= (z_{\mathrm{S},1} + z_\mathrm{P}) \omega_{8} - z_{\mathrm{P},1} \omega_{9}\\
z_{\mathrm{H},1} \omega_{4}  &= (z_{\mathrm{H},1} - z_\mathrm{P}) \omega_{8} + z_{\mathrm{P},1} \omega_{9} \\
\mathrm{Planetensatz\ 2}: \quad z_{\mathrm{S},2} \omega_{5}  &= (z_{\mathrm{S},2} + z_\mathrm{P}) \omega_{4} - z_{\mathrm{P},2} \omega_{10} \\
z_{\mathrm{H},2} \omega_{3}  &= (z_{\mathrm{H},2} - z_\mathrm{P}) \omega_{4} + z_{\mathrm{P},2} \omega_{10} \\
\mathrm{Planetensatz\ 3}: \quad z_{\mathrm{S},3} \omega_{3}  &= (z_{\mathrm{S},3} + z_\mathrm{P}) \omega_{2} - z_{\mathrm{P},3} \omega_{11} \\
z_{\mathrm{H},3} \omega_{6}  &= (z_{\mathrm{H},3} - z_\mathrm{P}) \omega_{2} + z_{\mathrm{P},3} \omega_{11} \\
\mathrm{Planetensatz\ 4}: \quad z_{\mathrm{S},4} \omega_{3}  &= (z_{\mathrm{S},4} + z_\mathrm{P}) \omega_{1} - z_{\mathrm{P},4} \omega_{12} \\
z_{\mathrm{H},4} \omega_{7}  &= (z_{\mathrm{H},4} - z_\mathrm{P}) \omega_{1} + z_{\mathrm{P},4} \omega_{12}.
	\end{align}
\end{subequations} 
\\
Die Gangschaltung erfolgt im Getriebe über jeweils drei Lamellenkupplungen und Lamellenbremsen. Mit Hilfe der Kupplungen lassen sich die Wellen 1 und 8, 8 und 3 und 2 und 7 verkoppeln. Über die drei Bremsen können die Wellen 8, 5 und 6 gegen das Gehäuse abgestützt werden. Um einen Gang zu schalten müssen drei der Schaltelemente betätigt sein. Dadurch verbleibt nur noch ein FHG der Getriebedynamik und es ergibt sich eine eindeutige Übersetzung zwischen Eingangsdrehzahl $\omega_{1}$ und Ausgangsdrehzahl $\omega_2$. Die sich durch die jeweiligen Schaltungen ergebenden Übersetzungen können Tabelle ??? entnommen werden.
\begin{table}
	\caption{Übersetzung und aktive Schaltelemente je Gang.}
	\label{tbl:clutches}
	\centering
	\begin{tabular}{l*{7}{c}} 
		\toprule
		& & \multicolumn{6}{c}{Schaltelement} \\
		\cmidrule{3-8}
		Gang & Übersetzung & K81 & B05 & B08 & K38 & K27 & B06 \\
		\midrule 
		1. Gang & 5,354 &           & $\bullet$ &           & $\bullet$ &           & $\bullet$ \\
		2. Gang & 3,243 & $\bullet$ &           &           & $\bullet$ &           & $\bullet$ \\
		3. Gang & 2,252 & $\bullet$ & $\bullet$ &           &           &           & $\bullet$ \\
		4. Gang & 1,636 &           & $\bullet$ &           &           & $\bullet$ & $\bullet$ \\
		5. Gang & 1,211 & $\bullet$ & $\bullet$ &           &           & $\bullet$ &           \\
		6. Gang & 1,000 & $\bullet$ &           &           & $\bullet$ & $\bullet$ &           \\
		7. Gang & 0,865 &           & $\bullet$ &           & $\bullet$ & $\bullet$ &           \\
		8. Gang & 0,717 &           &           & $\bullet$ & $\bullet$ & $\bullet$ &           \\
		9. Gang & 0,601 &           & $\bullet$ & $\bullet$ &           & $\bullet$ &           \\
		N       & $-$   &           & $\bullet$ &           &           &           & $\bullet$ \\
		R       & -4,798&           & $\bullet$ & $\bullet$ &           &           & $\bullet$ \\
		\bottomrule 
	\end{tabular} 
\end{table}

\subsection{Anwendung des Newton-Euler-Formalismus}\label{ssec:AnwNE}
Die Wellen und Planeten rotieren alle nur um eine Achse bzw. bewegen sich nur auf einer Kreisbahn. Daher werden die Orientierungen und Ortsvektoren des einzelnen Körper als skalare angegeben. Die Ortsvektoren werden als Zylinderkoordinaten geschrieben. Somit ergibt sich der Vektor der Orientierungen der einzelnen Wellen und der Planeten zu 
\begin{equation}
\pmb{s} = \begin{bmatrix} s_{1}\quad s_{2}\quad \dots \quad s_{12} \end{bmatrix}^T = \begin{bmatrix} \phi_{1}\quad \phi_{2}\quad \dots \quad \phi_{12} \end{bmatrix}^T.
\end{equation} 
Die Ortsvektoren der Wellen 1-8 können aufgrund der nicht vorhanden translatorischen Bewegung gleich Null gesetzt werden. Lediglich die Planeten haben durch die Bewegung auf den Planetenträgern um die Drehachse der Wellen~ 1-8 einen translatorischen Anteil. Die Position der Planeten ergibt sich somit aus den entsprechenden Radien der Planetenträgern $r_{T,j}$ und der Orientierung der Planetenträgern. Daraus ergibt sich der Orstvektor des Systems zu
\begin{equation}
\pmb{r} = \begin{bmatrix} r_{1}\; r_{2}\; \dots \; r_{12} \end{bmatrix}^T = \begin{bmatrix} 0\; \dots \; 0 \quad r_{T,1}\,\phi_8 \quad r_{T,2}\,\phi_4 \quad r_{T,3}\,\phi_2 \quad r_{T,4}\,\phi_1 \end{bmatrix}^T
\end{equation}
Da das MKS $f=4$ FHG hat lässt sich nach \cite{Schiehlen.2017} dieses durch verallgemeinerte Koordinaten $\pmb{y}\in \mathbb{R}^f$ beschreiben. Dabei liegt es nahe, die zu den am Getriebe gemessenen Winkelgeschwindigkeiten gehörigen Winkel $\phi_{1}$, $\phi_{2}$ und $\phi_{8}$ zu verwenden. Zusätzlich soll dazu $\phi_{3}$ verwendet werden. Dadurch ergibt sich der Vektor der verallgemeinerten Koordinaten zu
\begin{equation}
\pmb{y} = \begin{bmatrix} \phi_{1}\quad \phi_{2}\quad \phi_{3}\quad \phi_{8} \end{bmatrix}^T.
\end{equation}
Die Zusammenhänge in (\ref{eq:zwangsbedingungen}) gelten auch für die zugehörigen Winkel. Damit können $\pmb{s}$ und $\pmb{r}$ in Abhängigkeit der verallgemeinerten Koordinaten angegeben werden. Die Einträge von $\pmb{s}(\pmb{y})$ ergeben sich zu 
\begin{subequations}\label{eq:zwangsbedingungen_aufgelöst}
	\begin{align}
	 	\phi_1 &= \phi_1 \\
	 	\phi_2 &= \phi_2 \\
	 	\phi_3 &= \phi_3 \\ \label{eq:s4}
		\phi_{4} &= -\frac{z_{\mathrm{S,1}} }{z_{\mathrm{H,1}} } \phi_1 +\left(\frac{z_{\mathrm{S,1}} }{z_{\mathrm{H,1}} }+1\right) \phi_8 \\
		\phi_{5} &= \left(-\frac{z_\mathrm{S,1} }{z_\mathrm{H,1} }-\frac{z_\mathrm{H,2} z_\mathrm{S,1}}{z_\mathrm{H,1} z_\mathrm{S,2}} \right) \phi_1
			-\frac{z_\mathrm{H,2} }{z_\mathrm{S,2}} \phi_3
			+\left(\frac{z_{\mathrm{S,1}} }{z_{\mathrm{H,1}} }+\frac{z_{\mathrm{H,2}} }{z_{\mathrm{S,2}} }+\frac{z_{\mathrm{H,2}} z_{\mathrm{S,1}}}{z_{\mathrm{H,1}} z_{\mathrm{S,2}}} +1\right) \phi_8 \\
		\phi_{6} &= \left(\frac{z_{\mathrm{S,3}} }{z_{\mathrm{H,3}} }+1\right) \phi_2 - \frac{z_{\mathrm{S,3}} }{z_{\mathrm{H,3}} } \phi_3 \\
		\phi_{7} &= \left(\frac{z_{\mathrm{S,4}} }{z_{\mathrm{H,4}} }+1\right) \phi_1 - \frac{z_{\mathrm{S,4}} }{z_{\mathrm{H,4}} } \phi_3 \\
		\phi_8 &= \phi_8 \\
		\phi_{9} &= -\frac{z_{\mathrm{S,1}} }{z_\mathrm{P,1} } \phi_1 +\left(\frac{z_\mathrm{S,1} }{z_\mathrm{P,1}} +1\right) \phi_8 \\
		\phi_{10} &= \left(\frac{z_\mathrm{H,2} z_\mathrm{S,1}}{z_\mathrm{H,1} z_\mathrm{P,2}} + \frac{z_\mathrm{S,1}}{z_\mathrm{H,1}} \right) \omega_1
			+\frac{z_\mathrm{H,2}}{z_\mathrm{P,2}} \phi_3
			+\left(\frac{z_\mathrm{S,1}}{z_\mathrm{H,1}} -\frac{z_\mathrm{H,2} }{z_\mathrm{P,2} }-\frac{z_\mathrm{H,2}  z_\mathrm{S,1}}{z_\mathrm{H,1}  z_\mathrm{P,2}} + 1\right) \phi_8 \\
		\phi_{11} &= \left( \frac{z_{\mathrm{S,3}}}{z_{\mathrm{P,3}}} + 1 \right) \phi_2 -\frac{z_{\mathrm{S,3}}}{z_{\mathrm{P,3}}} \phi_3 \\
		\phi_{12} &= \left( \frac{z_{\mathrm{S,4}}}{z_{\mathrm{P,4}}} + 1 \right) \phi_1 -\frac{z_{\mathrm{S,4}}}{z_{\mathrm{P,4}}} \phi_3
	\end{align}
\end{subequations}
Die Einträge von $\pmb{r}$ sind bis auf den Zehnten bereits in Abhängigkeit von $\pmb{y}$ angegeben. Der verbleibende Eintrag kann mit (\ref{eq:s4}) wie folgt angegeben werden
\begin{equation}
r_{10}(\pmb{y}) = r_{T,2}\,\phi_4 = r_{T,2}\,\left(-\frac{z_{\mathrm{S,1}} }{z_{\mathrm{H,1}} } \phi_1 +\left(\frac{z_{\mathrm{S,1}} }{z_{\mathrm{H,1}} }+1\right) \phi_8\right).
\end{equation}
Durch die zeitliche Ableitung von $\pmb{r}(\pmb{y})$ und $\pmb{s}(\pmb{y})$ gemäß (\ref{vi}) und \ref{omegai} erhält man die Jacobi-Matrizen der Translation
\begin{equation}
\pmb{J}_\mathrm{T} = \begin{bmatrix}
J_\mathrm{T1} & J_\mathrm{T2} & \cdots & J_\mathrm{T12}
\end{bmatrix}^T
\end{equation}  
und der Rotation 
\begin{equation}
\pmb{J}_\mathrm{R} = \begin{bmatrix}
J_\mathrm{R1} & J_\mathrm{R2} & \cdots & J_\mathrm{R12}
\end{bmatrix}^T,
\end{equation} 
welche die einzelnen Jacobi-Matrizen der Körper enthalten. Die Terme $\overline{\pmb{v}}$ und $\overline{\pmb{\omega}}$ verschwinden, da es sich um skleronome Bindungen handelt und die ZB (\ref{eq:zwangsbedingungen}) nicht explizit zeitabhängig sind. Auch die lokalen Beschleunigungsvektoren $\overline{\pmb{a}}$ und $\overline{\pmb{\omega}}$ verschwinden, da $\pmb{J}_\mathrm{T}$ und $\pmb{J}_\mathrm{R}$ zeitlich konstant sind. In der allgemeinen Bewegungsgleichung für holonome MKS (\ref{Bwg}) verschwinden somit die Zentrifugalkräfte in $\pmb{k}(\pmb{y},\dot{\pmb{y}},t)$. Auch die Corioliskräfte verschwinden, da für alle Körper das Inertialsystem als Referenzsystem verwendet wird.%??? siehe TechDyn Seite 78 %
 Im Vektor der verallgemeinerten Kräfte $\pmb{q}(\pmb{y},\dot{\pmb{y}},t)$ werden die eingeprägten Kräfte $\pmb{f}^e_i$ vernachlässigt da die translatorischen Bewegungen im Getriebe sehr klein gegenüber der rotatorischen Bewegungen sind. Somit reduziert sich (\ref{Bwg}) auf
\begin{equation}
\pmb{M}(\pmb{y},t)\cdot\ddot{\pmb{y}} = \pmb{q}(\pmb{y},\dot{\pmb{y}},t)
\end{equation}
wobei die Massenmatrix mit $\pmb{J}_\mathrm{T}$ und $\pmb{J}_\mathrm{R}$ jetzt angegeben werden kann als 
  \begin{equation}
\pmb{M}(\pmb{y},t) = \pmb{J}_\mathrm{T}^T\cdot\mathrm{diag}(m_i)\cdot\pmb{J}_\mathrm{T}+\pmb{J}_\mathrm{R}^T\cdot\mathrm{diag}(I_i)\cdot\pmb{J}_\mathrm{R}
\end{equation}
und der Vektor der verallgemeinerten Kräfte als
  \begin{equation}
\pmb{q}(\pmb{y},\dot{\pmb{y}},t) = \pmb{J}_\mathrm{R}^T\cdot\mathrm{\pmb{l}}^e 
\end{equation}
Die eingeprägten Momente sind die von den Schaltelementen aufgebrachten Momente $\mathrm{T}_{\mathrm{B05}}$, $\mathrm{T}_{\mathrm{B06}}$, $\mathrm{T}_{\mathrm{B08}}$, $\mathrm{T}_{\mathrm{K81}}$, $\mathrm{T}_{\mathrm{K38}}$ und $\mathrm{T}_{\mathrm{K27}}$. Desweitern wirken am Getriebeeingang $\mathrm{T}_{\mathrm{In}}$ und am Getriebeausgang $\mathrm{T}_{\mathrm{Out}}$. Damit ergibt sich der Vektor der eingeprägten Momente zu
\begin{equation}\label{eq:le}
\mathrm{\pmb{l}}^e = \begin{bmatrix} \mathrm{T}_{\mathrm{In}}+\mathrm{T}_{\mathrm{K81}} \\ -\mathrm{T}_{\mathrm{K27}}-\mathrm{T}_{\mathrm{Out}} \\ -\mathrm{T}_{\mathrm{K38}} \\ 0 \\ \mathrm{T}_{\mathrm{B05}} \\ \mathrm{T}_{\mathrm{B06}} \\ \mathrm{T}_{\mathrm{K27}} \\ \mathrm{T}_{\mathrm{B08}}+\mathrm{T}_{\mathrm{K38}}-\mathrm{T}_{\mathrm{K81}} \\ 0 \\ 0 \\ 0 \\ 0 \end{bmatrix}.
\end{equation}
Das Eingangsmoment $\mathrm{T}_{\mathrm{In}}$ wird vom Verbrennungsmotor und einem Elektromotor gestellt. Das Ausgangssmoment $\mathrm{T}_{\mathrm{Out}}$ wird an die Kardanwelle übertragen. Andere eingeprägte Momente, wie Schleppmomente oder Reibungen an den Lagern werden vernachlässigt. Mit Hilfe der ersten Summe von (\ref{Bwg_lang}) können die Einträge der 4x4-Massenmatrix $\pmb{M}$ berechnet werden. Diese ergeben sich zu
\begin{align*}
	\pmb{M}_{(1,1)} &=  J_1+J_7\left(\frac{z_\mathrm{S,4}}{z_\mathrm{H,4}}+1\right)^2+J_{12} \left(\frac{z_\mathrm{S,4}}{z_\mathrm{P,4}}+1\right)^2+J_{10} {\left(\frac{z_\mathrm{S,1}}{z_\mathrm{H,1}} - \sigma_{5}\right)}^2+J_{5} \left(\frac{z_\mathrm{S,1}}{z_\mathrm{H,1}}+\sigma_{4}\right)^2 \\
		&\qquad + m_\mathrm{P,4} {r_\mathrm{T,4}}^2+\frac{J_{4} {z_\mathrm{S,1}}^2}{{z_\mathrm{H,1}}^2}+\frac{J_{9} {z_\mathrm{S,1}}^2}{{z_\mathrm{P,1}}^2}+\frac{m_\mathrm{P,2} {r_\mathrm{T,2}}^2 {z_\mathrm{S,1}}^2}{{z_\mathrm{H,1}}^2}\\
	%
	\pmb{M}_{(1,3)} &= \frac{J_{5} z_\mathrm{H,2} \left(\frac{z_\mathrm{S,1}}{z_\mathrm{H,1}}+\sigma_{4}\right)}{z_\mathrm{S,2}}-\frac{J_{12} z_\mathrm{S,4} \left(\frac{z_\mathrm{S,4}}{z_\mathrm{P,4}}+1\right)}{z_\mathrm{P,4}}-\frac{J_{10} z_\mathrm{H,2} \left(\frac{z_\mathrm{S,1}}{z_\mathrm{H,1}} - \sigma_{5}\right)}{z_\mathrm{P,2}}-\frac{J_{7} z_\mathrm{S,4} \left(\frac{z_\mathrm{S,4}}{z_\mathrm{H,4}}+1\right)}{z_\mathrm{H,4}}\\
	%
	\pmb{M}_{(1,4)} &= J_{10} \left(\frac{z_\mathrm{S,1}}{z_\mathrm{H,1}}-\sigma_{5}\right) \sigma_{3} - J_{5} \left(\frac{z_\mathrm{S,1}}{z_\mathrm{H,1}}+\sigma_{4}\right) \sigma_{2} -\frac{J_{4} z_\mathrm{S,1} \left(\frac{z_\mathrm{S,1}}{z_\mathrm{H,1}}+1\right)}{z_\mathrm{H,1}}\\
		&\qquad-\frac{J_{9} z_\mathrm{S,1} \left(\frac{z_\mathrm{S,1}}{z_\mathrm{P,1}}+1\right)}{z_\mathrm{P,1}}-\frac{m_\mathrm{P,2} {r_\mathrm{T,2}}^2 z_\mathrm{S,1} \left(\frac{z_\mathrm{S,1}}{z_\mathrm{H,1}}+1\right)}{z_\mathrm{H,1}}\\
	%
	\pmb{M}_{(2,2)} &= J_2+J_6 \left(\frac{z_\mathrm{S,3}}{z_\mathrm{H,3}}+1\right)^2+J_{11} \left(\frac{z_\mathrm{S,3}}{z_\mathrm{P,3}}+1\right)^2+m_\mathrm{P,3} {r_\mathrm{T,3}}^2\\
	%
	\pmb{M}_{(2,3)} &= -\frac{J_{6} z_\mathrm{S,3} \left(\frac{z_\mathrm{S,3}}{z_\mathrm{H,3}}+1\right)}{z_\mathrm{H,3}}-\frac{J_{11} z_\mathrm{S,3} \left(\frac{z_\mathrm{S,3}}{z_\mathrm{P,3}}+1\right)}{z_\mathrm{P,3}}\\
	\pmb{M}_{(3,1)} &= \pmb{M}_{(1,3)}\\
	\pmb{M}_{(3,2)} &= \pmb{M}_{(2,3)}\\
	\pmb{M}_{(3,3)} &= J_{3}+\frac{J_{10} {z_\mathrm{H,2}}^2}{{z_\mathrm{P,2}}^2}+\frac{J_{5} {z_\mathrm{H,2}}^2}{{z_\mathrm{S,2}}^2}+\frac{J_{6} {z_\mathrm{S,3}}^2}{{z_\mathrm{H,3}}^2}+\frac{J_{7} {z_\mathrm{S,4}}^2}{{z_\mathrm{H,4}}^2}+\frac{J_{11} {z_\mathrm{S,3}}^2}{{z_\mathrm{P,3}}^2}+\frac{J_{12} {z_\mathrm{S,4}}^2}{{z_\mathrm{P,4}}^2} \\
	%
	\pmb{M}_{(3,4)} &= -\frac{J_{10} z_\mathrm{H,2} \sigma_{3}}{z_\mathrm{P,2}} - \frac{J_{5} z_\mathrm{H,2} \sigma_{2}}{z_\mathrm{S,2}}\\
	\pmb{M}_{(4,1)} &= \pmb{M}_{(1,4)}\\
	\pmb{M}_{(4,3)} &= \pmb{M}_{(3,4)}\\
	%
	\pmb{M}_{(4,4)} &= J_8+J_4 \left(\frac{z_\mathrm{S,1}}{z_\mathrm{H,1}}+1\right)^2+J_{9} \left(\frac{z_\mathrm{S,1}}{z_\mathrm{P,1}}+1\right)^2+m_\mathrm{P,1} {r_\mathrm{T,1}}^2+J_{10} \sigma_{3}^2+J_5 \sigma_{2}^2 \\
		&\qquad + m_\mathrm{P,2} {r_\mathrm{T,2}}^2 \left(\frac{z_\mathrm{S,1}}{z_\mathrm{H,1}}+1\right)^2\\
	%
	\sigma_1 &= \left(\frac{z_{\mathrm{S,4}}}{z_{\mathrm{H,4}}}+1\right) \\
	\sigma_{2} &= \left(\frac{z_\mathrm{S,1}}{z_\mathrm{H,1}}+\frac{z_\mathrm{H,2}}{z_\mathrm{S,2}}+\frac{z_\mathrm{H,2} z_\mathrm{S,1}}{z_\mathrm{H,1} z_\mathrm{S,2}}+1\right)\\
	\sigma_{3} &= \left(\frac{z_\mathrm{H,2}}{z_\mathrm{P,2}}-\frac{z_\mathrm{S,1}}{z_\mathrm{H,1}}+\frac{z_\mathrm{H,2} z_\mathrm{S,1}}{z_\mathrm{H,1} z_\mathrm{P,2}}-1\right) \\
	\sigma_{5} &= \frac{z_\mathrm{H,2} z_\mathrm{S,1}}{z_\mathrm{H,1} z_\mathrm{P,2}} \\
	\sigma_{4} &= \frac{z_\mathrm{H,2} z_\mathrm{S,1}}{z_\mathrm{H,1} z_\mathrm{S,2}}
\end{align*}
wobei $J_{1-12}$ die Massenträgheitsmomente der entsprechenden Wellen sind und $m_{\mathrm{P},1-4}$ die gesamt Masse der Planeten eines Planetensatzes. Durch einsetzten von $\pmb{M}$ und \eqref{eq:le} in \eqref{Bwg_lang} ergeben sich die vier Gleichungen der Getriebedynamik
\begin{subequations} \label{eq:bewegungs_gleichungen}
	\begin{align}
		\pmb{M}_{(1,1)} \dot{\omega}_1 + \Vec{M}_{(1,3)} \dot{\omega}_3 +\pmb{M}_{(1,4)} \dot{\omega}_8
		&= T_{\mathrm{K81}} +T_{\mathrm{in}} + \sigma_1 T_{\mathrm{K27}}-T_{\mathrm{B05}} {\left(\frac{z_{\mathrm{S,1}} }{z_{\mathrm{H,1}}}+\sigma_{4} \right)}\\
		%
		\pmb{M}_{(2,2)} \dot{\omega}_2 + \pmb{M}_{(2,3)} \dot{\omega}_3
		&= T_{\mathrm{B06}} {\left(\frac{z_{\mathrm{S,3}} }{z_{\mathrm{H,3}} }+1\right)}-T_{\mathrm{out}} -T_{\mathrm{K27}} \\
		%
		\pmb{M}_{(3,1)} \dot{\omega}_1 + \pmb{M}_{(3,2)} \dot{\omega}_2 + \pmb{M}_{(3,3)} \dot{\omega}_3 + \pmb{M}_{(3,4)} \dot{\omega}_8
		&= -T_{\mathrm{K38}} -\frac{ z_{\mathrm{H,2}} }{z_{\mathrm{S,2}} }T_{\mathrm{B05}}-\frac{z_{\mathrm{S,3}}}{z_{\mathrm{H,3}}}T_{\mathrm{B06}}-\frac{z_{\mathrm{S,4}} }{z_{\mathrm{H,4}}}T_{\mathrm{K27}}\\
		%
		\pmb{M}_{(4,1)} \dot{\omega}_1 + \pmb{M}_{(4,3)} \dot{\omega}_3 + \pmb{M}_{(4,4)} \dot{\omega}_8
		&= T_{\mathrm{B08}} +T_{\mathrm{K38}} -T_{\mathrm{K81}} +\sigma_{2} T_{\mathrm{B05}}.
	\end{align}
\end{subequations}
Mit Hilfe der Definition des Eingangsvektors
\begin{equation}\label{eq:u}
\pmb{u} = \begin{bmatrix} \mathrm{T}_{\mathrm{In}} \quad \mathrm{T}_{\mathrm{Out}} \quad \mathrm{T}_{\mathrm{K81}} \quad \mathrm{T}_{\mathrm{K38}} \quad \mathrm{T}_{\mathrm{B08}} \quad \mathrm{T}_{\mathrm{B05}} \quad \mathrm{T}_{\mathrm{B06}} \quad \mathrm{T}_{\mathrm{K27}} \end{bmatrix}^T
\end{equation}
lassen sich die Gleichungen \eqref{eq:bewegungs_gleichungen} schreiben als
\begin{equation}
\pmb{M}\,\ddot{\pmb{y}} = \tilde{\pmb{B}}\,\pmb{u}.
\end{equation}
Dieses Differentialgleichungssystem kann mit der Definition des Zustands
\begin{equation}
\pmb{x} = \dot{\pmb{y}} = \begin{bmatrix} \omega_1 \\ \omega_2 \\ \omega_3 \\ \omega_8\end{bmatrix}
\end{equation}
und der Invertierung von $\pmb{M}$ in die Zustandsraumdarstellung 
\begin{equation}
\dot{\pmb{x}} = \pmb{M}^{-1}\tilde{\pmb{B}}\,\pmb{u} = \pmb{B}\,\pmb{u}
\end{equation}
transformiert werden. In dieser Darstellung ist leicht ersichtlich, dass das System aufgrund der fehlenden Systemmatrix zustandsunabhängig ist.

\subsection{Erweiterung um elastische Seitenwellen}
Aufgrund der relativ elastischen Seitenwellen kann es im Antriebsstrang zu erheblichen Schwingungen kommen. Um diese bei der Simulation eines Schaltvorgang mit berücksichtigen zu können, soll das in \ref{ssec:AnwNE} hergeleitete Zustandsraummodell um die Elastizitäten erweitert werden. Die Seitenwellen werden dabei als Feder-Dämpfer-Systeme modelliert. Die Kardanwelle wird hingegen als starrer Körper betrachtet. Ein schematischen Aufbau des Antriebsstrangs zeigt Abbildung ???. Zunächst wird der Zustandsvektor um die Verdrehung der Seitenwellen~$\psi$ und die Raddrehzahl~$\omega_\mathrm{C}$ ergänzt, sodass sich der erweiterte Zustandsvektor
\begin{equation}
\pmb{x}_\mathrm{ex} = \begin{bmatrix} \omega_1 \\ \omega_2 \\ \omega_3 \\ \omega_8 \\ \phi \\ \omega_\mathrm{C} \end{bmatrix}
\end{equation}
ergibt. Wie in Abbildung ??? veranschaulicht, ergibt sich die $\phi$ aus der Differenz der Verdrehung am linken Schnittufer der Seitenweillen $\phi_\mathrm{ss,p}$ und der des Rades $\phi_\mathrm{C}$, am rechten Schnittufer. Unter Berücksichtigung der starren Kardanwelle ergibt sich somit für die Verdrehung der Seitenwellen
\begin{equation}
\phi = \frac{\phi_2}{i_{\mathrm{D}}} - \phi_\mathrm{C}
\end{equation}  
wobei $i_{\mathrm{D}}$ die Übersetzungen des Differentials angibt. Entsprechend gilt der Zusammenhang auch für die Winkelgeschwindigkeiten
\begin{equation}\label{eq:dynphi}
\dot{\phi} = \frac{\omega_2}{i_{\mathrm{D}}} - \omega_\mathrm{C}.
\end{equation}
Des weiteren kann am Rad das Momentengleichgewicht
\begin{equation}\label{eq:MggRad}
I^\mathrm{eff}_\mathrm{C}\,\dot{\omega}_\mathrm{C} = 2\,T_\mathrm{ss} - T_\mathrm{res}
\end{equation}
aufgestellt werden. Das effektive Massenträgheitsmomentes des Rads $I^\mathrm{eff}_\mathrm{C}$ setzt sich wie folgt zusammen
\begin{equation}
I^\mathrm{eff}_\mathrm{C} = I_\mathrm{ss} + m_\mathrm{C}\,r_\mathrm{dyn}.
\end{equation}
Dabei ist $I_\mathrm{ss}$ das Massenträgheitsmoment einer Seitenwelle, $m_\mathrm{C}$ die Masse des Rads und $r_\mathrm{dyn}$ der dynamische Radradius. Desweiteren ist $T_\mathrm{ss}$ das von den Seitenwellen übertragene Moment und kann mit dem Steifigkeitskoeffizient $k_\mathrm{ss}$ und dem Dämpfungskoeffizient $d_\mathrm{ss}$ berechnet werden zu
\begin{equation}\label{eq:Tss}
T_\mathrm{ss} = k_\mathrm{ss}\left(\frac{\phi_2}{i_{\mathrm{D}}} - \phi_\mathrm{C}\right) + d_\mathrm{ss}\left(\frac{\omega_2}{i_{\mathrm{D}}} - \omega_\mathrm{C}\right).
\end{equation}
Das Widerstandsmoment berechnet sich als Summe der berücksichtigten äußeren Widerstände zu
\begin{equation}\label{eq:Tres}
T_\mathrm{res} = r_\mathrm{dyn}\left(F_\mathrm{r}(v) + F_\mathrm{ad}(v) + F_\mathrm{g}\right). 
\end{equation}
Die berücksichtigten Widerstände sind der Rollwiderstand $F_\mathrm{r}(v)$, der Luftwiderstand $F_\mathrm{ad}(v)$ und die Hangabtriebskraft $F_\mathrm{g}$ und werden wie folgt berechnet
\begin{subequations}
\begin{align}
F_\mathrm{r}(v) &= v \cdot f_\mathrm{Roll}(v)\cdot \cos\left(\zeta(t)\right)\cdot m_\mathrm{veh}\cdot g \\
F_\mathrm{ad}(v) &= \frac{1}{2}\cdot v^2\cdot\rho_\mathrm{Air}\cdot c_\mathrm{W}\cdot A_\mathrm{front}\\
F_\mathrm{g} &= \sin\left(\zeta(t)\right)\cdot m_\mathrm{veh}\cdot g.
\end{align}
\end{subequations}
Dabei ist $v = \omega_\mathrm{C}\,r_\mathrm{dyn}$ die Fahrzeuggeschwindigkeit, $f_\mathrm{Roll}(v)$ der von der Geschwindigkeit abhängige Faktor
des Rollwiderstands, $\zeta$ die Steigung der Straße in rad, $g$ die Erdbeschleunigung, $\rho_\mathrm{Air}$ die Luftdichte, $c_\mathrm{W}$ der Widerstandsbeiwert des Fahrzeugs, $A_\mathrm{front}$ die Frontfläche des
Fahrzeugs und $m_\mathrm{veh}$ die Gesamtmasse des Fahrzeugs \cite[S.~3ff]{Naunheimer.2007}.
Durch Einsetzten von \eqref{eq:Tss} und \eqref{eq:Tres} in \eqref{eq:MggRad} und anschließender Auflösung nach $\dot{\omega}_\mathrm{C}$ erhält man die Dynamik des Rades
\begin{equation}\label{eq:dynwc}
\dot{\omega}_\mathrm{C} = \left[ 2\,\left(k_\mathrm{ss}\left(\frac{\phi_2}{i_{\mathrm{D}}} - \phi_\mathrm{C}\right) + d_\mathrm{ss}\left(\frac{\omega_2}{i_{\mathrm{D}}} - \omega_\mathrm{C}\right)\right) - r_\mathrm{dyn}\left(F_\mathrm{r}(v) + F_\mathrm{ad}(v) + F_\mathrm{g}\right)\right]/I^\mathrm{eff}_\mathrm{C}.
\end{equation}
Aufgrund der starren Verbindung von Welle 2, der Kardanwelle und des Differentials, werden die Massenträgheitsmomente dieser Bauteile zusammengefasst zu
\begin{equation}
J^\mathrm{eff}_2 = J_2 + J^\mathrm{eff}_\mathrm{cs} + \frac{J^\mathrm{eff}_\mathrm{cs}}{i_\mathrm{D}^2}
\end{equation}
wobei in das effektive Massenträgheitsmoment der Kardanwelle $J^\mathrm{eff}_\mathrm{cs} = J_\mathrm{cs} + 0,5\,J_\mathrm{D}$ die eine Hälfte des Massenträgheitsmoments des Differentials miteinbezogen wird. Die andere Hälfte wird aufgrund der Drehzahldifferenz nach dem Differential mit dem Massenträgheitsmoment der zweiten Seitenwelle $J^\mathrm{eff}_\mathrm{ss} =  0,5\,J_\mathrm{D} + J_\mathrm{ss}$ zusammen gefasst. Für das erweiterte System wird $J_2$ in $\pmb{M}_(2,2)$ dann durch $J^\mathrm{eff}_2$ ersetzt. Des weiteren wird in Abbildung ??? ersichtlich, dass für das Ausgangsmoment 
\begin{equation}
T_\mathrm{out} = \frac{2\,T_\mathrm{ss}}{i_\mathrm{D}} = \frac{2}{i_\mathrm{D}}\,\left[ k_\mathrm{ss}\left(\frac{\phi_2}{i_{\mathrm{D}}} - \phi_\mathrm{C}\right) + d_\mathrm{ss}\left(\frac{\omega_2}{i_{\mathrm{D}}} - \omega_\mathrm{C}\right)\right]
\end{equation}
gilt und dieses in \eqref{eq:le} dementsprechend ersetzt werden kann.

Zu beachten ist hier, dass $v(\omega_\mathrm{C})$ zustandsabhängig ist und sowohl in $F_\mathrm{r}(v)$ als auch in $F_\mathrm{ad}(v)$ nichtlinear eingeht, wodurch das gesamt System nichtlinear wird. Auch die Eingangsgröße $\zeta(t)$ geht nichtlinear in die Widerstand ein. Da Fahrbahnen im Straßenverkehr auf maximal $20^\circ$ beschränkt sind werden hier die Vereinfachungen
\begin{equation}
\sin(\zeta) = \zeta \quad \mathrm{und} \quad \cos(\zeta) = 1
\end{equation}  
getroffen. Damit lässt sich das erweiterte Modell mit dem ergänzten Eingangsvektor
\begin{equation}\label{eq:uex}
\pmb{u}_\mathrm{ex,nl} = \begin{bmatrix} \mathrm{T}_{\mathrm{In}} & \mathrm{T}_{\mathrm{K81}} & \mathrm{T}_{\mathrm{K38}} & \mathrm{T}_{\mathrm{B08}} & \mathrm{T}_{\mathrm{B05}} & \mathrm{T}_{\mathrm{B06}} & \mathrm{T}_{\mathrm{K27}} & \zeta \end{bmatrix}^T,
\end{equation}
den Gleichungen \eqref{eq:dynphi} und \eqref{eq:dynwc} in der Form 
\begin{equation}
\dot{\pmb{x}}_\mathrm{ex} = \pmb{f}(\pmb{x}_\mathrm{ex}) + \pmb{B}_\mathrm{ex,nl}\,\pmb{u}_\mathrm{ex,nl}
\end{equation}
schreiben.

Des weiteren kann durch die Betrachtung der Widerstandskräfte $F_\mathrm{r}(v)$ und $F_\mathrm{ad}(v)$ als Eingangsgrößen, das System weiterhin in lineare Form
\begin{equation}\label{eq:sys_linex}
\dot{\pmb{x}}_\mathrm{ex} = \pmb{A}_\mathrm{ex}\,\pmb{x}_\mathrm{ex} + \pmb{B}_\mathrm{ex}\,\pmb{u}_\mathrm{ex}
\end{equation}
angegeben werden. Hierfür werden die beiden Widerstandskräfte in \eqref{eq:dynwc} durch $T^*_\mathrm{res}$ ersetzt, sodass sich 
\begin{equation}\label{eq:dynwc_lin}
\dot{\omega}_\mathrm{C} = \left[ 2\,\left(k_\mathrm{ss}\left(\frac{\phi_2}{i_{\mathrm{D}}} - \phi_\mathrm{C}\right) + d_\mathrm{ss}\left(\frac{\omega_2}{i_{\mathrm{D}}} - \omega_\mathrm{C}\right)\right) - T^*_\mathrm{res} - r_\mathrm{dyn}\, F_\mathrm{g}\right]/I^\mathrm{eff}_\mathrm{C}.
\end{equation}
 ergibt. Der Eingangsvektor des linearen Systems ist dann definiert als 
\begin{equation}\label{eq:uex}
\pmb{u}_\mathrm{ex} = \begin{bmatrix} \mathrm{T}_{\mathrm{In}} & T^*_\mathrm{res} & \mathrm{T}_{\mathrm{K81}} & \mathrm{T}_{\mathrm{K38}} & \mathrm{T}_{\mathrm{B08}} & \mathrm{T}_{\mathrm{B05}} & \mathrm{T}_{\mathrm{B06}} & \mathrm{T}_{\mathrm{K27}} & \zeta \end{bmatrix}^T.
\end{equation}

Die Funktion $\pmb{f}(\pmb{x}_\mathrm{ex})$ und die Eingangmatrix $\pmb{B}_\mathrm{ex,nl}$ des nichtlinearen Systems, sowie die Systemmatrix $\pmb{A}_\mathrm{ex}$ und die Eingangmatrix $\pmb{B}_\mathrm{ex}$ des linearen Systems sind im Anhang mit numerischen Werten für die Zähnezahlen der Zahnräder und der Massenträgheitsmomente angegeben.

\section{Modellierung der Reibung an Kupplungen und Bremsen}
Die im Abschnitt \ref{ssec:AnwNE} eingeführten Momente der Schaltelemente werden über Lamellenkupplungen beziehungsweise -bremsen an den Wellen angebracht. Dabei werden die Lamellenpakete über einen hydraulischen Druck zusammengepresst, sodass zwischen den Flächen Reibkräfte wirken. Solang die aus dem Druck resultierende Normalkraft an Lamellenpaketen nicht ausreicht um die Scheiben gegeneinander abzustützen, wirkt die Reibkraft entgegen der Differnzdrehzahl $\Delta \omega$ der drehenden Scheiben. Ist die Normalkraft groß genug entsteht eine kraftschlüssige Verbindung zwischen Innen- und Außenlamellen. In diesem Fall ist $\Delta \omega = 0$ und der Betrag der Reibkraft entspricht dem der Reaktionskraft und wirkt dieser entgegen.   

Zur Ermittlung der auf die Wellen wirkenden Reaktionsmomente bedarf es eines Reibmodells. Im Wesentlichen können diese in statische und dynamische Reibmodelle eingeteilt werden. Dynamische Reibmodelle unterscheiden sich dadurch, dass sie über mindestens einen weiteren Zustand verfügen. Dies ermöglicht die Simulation von weiteren Reibeffekten wie Hysteresen, variierende Losbrechkräfte oder nicht-gleitende Auslenkungen \cite{Schroeder.2015}. Im folgenden werden statische und dynamische Reibmodelle erläutert.

\subsection{Statische Reibmodelle}
Die einfachste Art der Reibungsmodellierung ist die Coulombsche Reibung. Dabei ist die Reibkraft proportional zur Normalkraft $F_\mathrm{N}$ und einem konstanten Reibungskoeffizienten $\mu_\mathrm{C}$. Nach \cite{Schroeder.2015} ist diese definiert als
\begin{equation}
F_\mathrm{C} = F_\mathrm{N}\,\mu_\mathrm{C}\;\mathrm{sign}(\Delta \omega)=F_\mathrm{N}\,\mu_\mathrm{C}\begin{cases}1&\quad,\Delta \omega>0\\0&\quad,\Delta \omega=0\\-1&\quad,\Delta \omega<0 \end{cases}
\end{equation}
Durch die Erweiterung um einen Haftreibungsbeiwert $\mu_\mathrm{S}$, wird eine benötigte Losbrechkraft mit berücksichtigt. Der Haftreibungsbeiwert $\mu_\mathrm{S}$ ist im Normalfall größer als $\mu_\mathrm{C}$, sodass eine größere Kraft notwendig ist um einen reibungsbehafteten Körper aus der ruhe zu bringen. Für die Haftreibung gilt
\begin{equation}
F_\mathrm{S} = \begin{cases}F_\mathrm{A}&,\vert F_\mathrm{A}\vert < F_\mathrm{N}\,\mu_\mathrm{S} (\wedge\Delta \omega = 0)
\\ F_\mathrm{N}\,\mu_\mathrm{S}\,\mathrm{sign}(F_\mathrm{A})&,\vert F_\mathrm{A}\vert \geq F_\mathrm{N}\,\mu_\mathrm{S} (\wedge\Delta \omega = 0) \end{cases}
\end{equation}
wobei $F_\mathrm{A}$ die Summe der äußeren Kräfte ist. Sind die Reibflächen geschmiert, ist eine Änderung der Reibkennlinie über  die $\Delta \omega$ zu beobachten. Dieses Phänomen berücksichtigt die viskose Reibung. Diese ist definiert als
\begin{equation}
F_\mathrm{V}=\nu\vert\Delta\omega\vert^{\delta_\mathrm{V}}\mathrm{sign}(\Delta\omega)
\end{equation}
mit dem viskosen Reibungsexponenten $\delta_\mathrm{V}$ und dem Reibkoeffizienten $\nu$. Des weiteren lässt sich $F_\mathrm{C}$ durch die Stribeck-Funktion erweitern. Damit wird das rapide Abfallen der Reibkraft bei kleinen Geschwindigkeiten berücksichtigt. Dieser Effekt entsteht durch das Aufbrechen von Kontaktstellen und der anschließenden Bildung eines Schmierfilms \cite{Stribeck.1903}. Eine übliche Modellierung der Stribeck-Funktion ist
\begin{equation}
F_\mathrm{Stribeck} = F_\mathrm{C} + F_\mathrm{N}(\mu_\mathrm{S}-\mu_\mathrm{C})\mathrm{exp}\left[-\left(\frac{\vert\Delta \omega \vert}{\omega_\mathrm{S}}\right)^{\delta_\mathrm{S}}\right]
\end{equation} 
mit der Stribeck-Geschwindigkeit $\omega_\mathrm{S}$ und dem Stribeck-Exponent $\delta_\mathrm{S}$. Die verschiedene statischen Reibmodelle lassen sich zum "Kinetischen Reibmodell" 
\begin{equation}
F_\mathrm{Kin} = \begin{cases}F_\mathrm{A} &,\Delta \omega=0 \\
 F_\mathrm{Stribeck}\; \mathrm{sign}(\Delta \omega)+F_\mathrm{V} &, \Delta \omega \neq 0\end{cases}
\end{equation}
zusammenfassen \cite{Altpeter.1999}.

Eine Schwäche dieses Modells ist der unstetige Verlauf der Reibkennlinie bei $\Delta \omega=0$. Diese Schwäche kann durch die Annäherung der Reibkennlinie über eine stetige atan-Funktion behoben werden. Dennoch bedarf es weiterhin der Unterscheidung zwischen Haft- und Gleitbereich. Hierfür ist entweder eine variable Zeitschrittweite des Solvers notwendig oder die Definition eines Geschwindigkeitbandes um $\Delta\omega=0$. Des weiteren ist die Reibkraft im Haftbereich von den äußeren Kräften abhängig, welche hier nicht bekannt sind und durch das Modell berechnet werden sollen. Aus diesen Gründen werden im Folgenden dynamische Reibmodelle betrachtet, welche diese Schwächen nicht aufweisen. 

\subsection{Dynamische Reibmodelle}
Um die dynamischen Reibeffekte vorallem um $\Delta\omega=0$ mit berücksichtigen zu können werden dynamische Reibmodelle benötigt. Einen ersten Ansatz dazu liefert das Reibmodell von Dahl. Dieses beschreibt die dynamischen Reibeffekte für kleine Auslenkungen mit kleinen Federn in den Kontaktbereichen von Körpern und geht anschließend in die Coulombschen-Reibung über\cite{Schroeder.2015}. Die Reibdynamik des Dahlschen-Modells wird definiert durch die Differentialgleichung
\begin{equation}
F_\mathrm{R,Dahl} = \left[\sigma_0\vert 1-\frac{F_\mathrm{R,Dahl}}{F_\mathrm{C}}\mathrm{sign}(\Delta\omega)\vert^\beta\ \mathrm{sign}\left( 1-\frac{F_\mathrm{R,Dahl}}{F_\mathrm{C}}\mathrm{sign}(\Delta\omega)\right)\right]\Delta\omega
\end{equation}
mit der zuvor definierten Coulomschen-Reibkraft $F_\mathrm{C}$, der Federsteifigkeit $\sigma_0$ und dem Formparameter der Hysterese $\beta$. Das Lund-Grenoble (LuGre) Reibmodell erweitert das Modell von Dahl um einen Dämpfungsanteil und ersetzt die Coulombsche-Kraft durch einen beliebige Reibkraftkennlinie $F_\mathrm{SS}$. Die Kontaktstellen der Reibflächen werden dabei als "Borsten"  modelliert, welche die Borstensteifigkeit $\sigma_0$ und die Borstendämpfung $sigma_1$ besitzen. Mit der Borstenauslenkung $z$ lässt sich die LuGre-Reibkraft angeben als
\begin{equation}
F_\mathrm{R,LuGre} = \sigma_0\, z + \sigma_1 \dot{z} + \sigma_2\, \Delta \omega\quad \mathrm{mit}\  \dot{z}=\Delta\omega\left(1-\frac{\sigma_0 z}{F_\mathrm{SS}}\right)
\end{equation} 
wobei $\sigma_2$ der Steigungskoeffizienten der viskosen Reibung ist [Quelle aus Mahl ???]. Jedoch wird in \cite{Dupont.2002} gezeigt, dass dieses Modell kein elasto-plastisches Presliding abbilden kann. Daher führt jede nichtkonstante Reibkraft zu einem sogenannten plastischen Presliding von $z$ wodurch es besonders bei periodischen Kräften zu einem physikalisch nicht existierenden Drift kommt. Das Elasto-Plastische Modell erweitert das LuGre-Modell um einen elastischen Bereich für kleine $z$. Nach \cite{Dupont.2000} ergibt sich die Dynamik der Auslenkung zu 
\begin{equation}
\dot{z} = \Delta \omega \left(1-\alpha(z,\Delta \omega)\frac{\sigma_0\,z}{F_\mathrm{SS}}\right)^i,\quad i\in \mathbb{Z}
\end{equation}
wobei $i=1$ gesetzt wird und die Funktion $\alpha(z,\Delta \omega)$ ein elastische Bereich für kleine $z$ erzeugt.

Aufgrund der erläuterten Vorteile wird das Elasto-Plastische-Modell für das Simulationsmodell verwendet. Mit dem mittleren Radius $r_\mathrm{m}$ der Kupplungs- und Bremsscheiben und der Anzahl der Scheiben pro Schaltelement $N_\mathrm{disc}$  lassen sich die oben beschrieben Kräfte in Momente umrechnen. Letztendlich ergibt sich das an der Welle angreifende Reibmoment des Schaltelements $k$ zu
\begin{align}
T_{\mathrm{Cl},k} = \sigma_0\,z + \sigma_1\,\dot{z} + \sigma_2 \Delta\omega
\end{align}
mit
\begin{align}
\dot{z} = \Delta \omega \left(1-\alpha\frac{\sigma_0\,z}{T_\mathrm{Cl,ss}}\right)\\
T_\mathrm{Cl,ss} = 2\,N_{\mathrm{Disc},k}\,r_{\mathrm{m},k}\,F_\mathrm{Stribeck}.
\end{align}
Die Reibwerte $\mu_\mathrm{S}$ und $\mu_\mathrm{C}$ werden aus Kennfeldern in Abhängigkeit der des Drucks an den Scheiben $p_\mathrm{N}$, der Öltemperatur $T_\mathrm{Oil}$ und $\Delta \omega$ gewonnen. Die Normalkraft $F_\mathrm{N}$ ergibt sich aus $p_\mathrm{N}$ und der Fläche der Scheiben $A_\mathrm{Disc}$ und wird vorgegeben. 
\chapter{Sensitivitätsanalyse}\label{ch:Sensitivity}\label{ch:ch3}
In diesem Kapitel wird das in Kapitel \cite{ch:ch3} hergeleitete lineare Modell des Antriebsstrangs auf die Sensitivität verschiedener Parameter und Eingänge untersucht. Damit soll die Relevanz für eine Parameter- und Eingangsschätzung im folgenden Kapitel belegt und unterstrichen werden.
Eine mögliche Definition für eine Sensitivitätsanalyse ist die folgende: \emph{The study of how uncertainty in the output of a model (numerical or otherwise) can be apportioned to different sources of uncertainty in the model input} \cite{Saltelli.2004}.
Für die Analyse gibt es verschiedene Methoden. Dabei ist die OAT(one-at-a-time)-Methode wohl die einfachste. Dabei wird pro Simulationsdurchlauf lediglich eine Eingangsgröße geändert während alle anderen auf den nominal Werten gehalten werden. Die Methode liefert aber nur eine sehr beschränkte Aussage über den gesamten Eingangsraum. Jedoch kann sie für die Abschätzung einer ersten Untergruppe an relevanten Eingängen hilfreich sein.    

 Die in der Literatur am häufigsten diskutierten Ansätze basieren auf partiellen Ableitungen der Art $\partial y_i/\partial x_i$ wobei $y$ ein beliebiger Ausgang und $x_i$ ein beliebiger Eingang des Systems ist. Ein Eingang kann hierbei auch ein variierender Parameter sein. Dieses Vorgehen wird vor allem bei relativ einfachen, mathematisch wohlbestimmten und stetigen Modellen verwendet. Der Vorteil dabei ist, dass die Sensitvitätsfunktionen sobald sie einmal berechnet sind leicht an Veränderungen anpassbar sind und auch für ähnliche Systeme verwendet werden können \cite{Karnavas.1993}. Jedoch werden die partiellen Ableitung immer in einem bestimmen Arbeitspunkt berechnet d.h. die Aussagen der Sensitvitätsfunktionen gelten für nichtlineare Systeme nur lokal um diesen den Arbeitspunkt. Des weiteren ist es oft sehr schwer oder gar unmöglich die benötigten Ableitungen für komplexe Modelle zu berechnen.

Eine Alternative zur Berechnung der partiellen Ableitungen sind globale Methoden. Hierbei werden die Modelleingänge stochstisch varriert, woraus sich nach vielen Simulationen eine Strichprobe für die Modellausgänge ergibt. Dieses Methode der mehrfachen Simulation wird auch Monte-Carlo-Methode genannt. Wie der Name schon sagt ist der grundlegende Vorteil dabei, dass durch eine genügend große Stichprobe eine globale Aussage zur Sensitivität gegeben werden kann. Des weiteren erlaubt die Methode auch Aussagen über die verkoppelte Wirkung von Eingangsänderungen ohne spezielles Vorwissen über die Struktur des Modells. Der entscheidende Nachteil liegt je nach Stichprobengröße und Komplexität des Modells in der Rechenzeit.      

Aufgrund des linearen und stetigen Charakters von \eqref{eq:sys_linex} werden für diese Arbeit die partiellen Ableitungen für die Sensitivitätsanalyse verwendet.
Hierfür wird zuerst ein idealer Schaltvorgang beschrieben um die Sensitivitätsfunktionen in verschieden Schaltphasen auswerten zukönnen. Des weiteren wird mit der Berechnung der Fisher-Informationsmatrix eine Abschätzung für die erreichbare Genauigkeit der Schätzungen geliefert.
   
\section{Beschreibung eines Gangwechsels}\label{sec:Gangwechsel}
Bei einem Gangwechsel wird durch die Betätigung von Schaltelementen die Übersetzung des Getriebes geändert. Dies soll zum einen in jedem Fall möglichst ohne für den Fahrer spürbare unerwünschte Beschleunigungen am Getriebeabtrieb passieren. Zum anderen soll der Gangwechsel aber auch so schnell wie möglich und ohne Momentenunterbrechung vollzogen werden. Es besteht also eine Zielkonflikt zwischen den Komfortansprüchen des Fahrers und der Dauer eines Schaltvorgangs. Grundsätzlich wird zwischen vier verschieden Schaltvorgängen unterschieden. 
Bei dem in dieser Arbeit betrachteten Getriebe geschieht ein Gangwechsel mittels der Schließung und der simultanen Öffnung einer Bremse bzw. Kupplung. In der   

\section{Anwendung der Paramater- und Eingangssensitivitätsanalyse}
Wie einleitend erläutert werden die Sensitiviäten in dieser Arbeit mittels partiellen Differentialen berechnet. Nach \cite{Hearne.1985} ist für ein nichtlineares System erster Ordnung
\begin{equation}
\mathrm{d} x_i/\mathrm{d} t = f_i(\pmb{x},\pmb{p},\pmb{u})\quad i=1,2,\,\dots\, ,n 
\end{equation}
mit dem Zustandsvektor $\pmb{x} = [x_1, x_2,\, \dots\, ,x_n]^T$, dem Parametervektor $\pmb{p} = [p_1, p_2,\, \dots\, ,p_m]^T$, dem Eingangsvektor $\pmb{u} = [u_1, u_2,\, \dots\, ,p_l]^T$ und der differenzierbaren Funktion $f_i$ die absolute Parametersensitivitätsfunktion definiert als
\begin{equation}\label{eq:abs_psens}
S_{ij} = \partial x_i/\partial p_j \quad i=1,2,\,\dots\, ,n \quad j=1,2,\,\dots\, ,m.
\end{equation}
Die entsprechende normalisierte Parametersensitivitätsfunktion ist 
\begin{equation}
S_{\mathrm{N},ij} = \partial x_i/\partial p_j\ p_{j0}/x_{i0}.
\end{equation}
Die normalisierte Parametersensitivitätsfunktion $S_{ij}$ kann interpretiert werden, als approximierte prozentuale Änderung des Zustandes $x_i$ für eine $1\%$-ige Änderung des Parameters $p_j$. Dem entsprechend sind auch die absolute Eingangssensitivitätsfunktion
\begin{equation}\label{eq:abs_usens}
S_{ik} = \partial x_i/\partial u_k \quad i=1,2,\,\dots\, ,n \quad k=1,2,\,\dots\, ,l
\end{equation}
und die normalisierte Eingangssensitivitätsfunktion
\begin{equation}\label{eq:abs_usens}
S_{ik} = \partial x_i/\partial u_k\ u_{k0}/x_{i0}
\end{equation}
definiert.
Wendet man \eqref{eq:abs_psens} auf das lineare System 
\begin{equation}
\mathrm{d} \pmb{x}/\mathrm{d} t = \pmb{A}(\pmb{p})\,\pmb{x}+\pmb{B}(\pmb{p})\,\pmb{u}
\end{equation}
an, erhält man
\begin{equation}
\mathrm{d} \pmb{S}_j/\mathrm{d} t = \dot{\pmb{S}}_j= \pmb{A}(\pmb{p})\,\pmb{S}_j+\partial/\partial p_j\ \pmb{B}(\pmb{p})\,\pmb{u}.
\end{equation}
Für die Eingangssensitivität ergibt sich 
\begin{equation}
\mathrm{d} \pmb{S}_k/\mathrm{d} t = \dot{\pmb{S}}_k= \partial/\partial u_i\ \pmb{B}(\pmb{p})\,\pmb{u}.
\end{equation}
Die Parameter- und Eingangsensitivitätsfunktion werden im Folgenden anhand \eqref{eq:sys_ex23lin} für $m_\mathrm{veh}$, $k_\mathrm{ss}$, $d_\mathrm{ss}$, $\mu_{38}$ und $\zeta$ gebildet und die zeitlichen Verläufe von $\pmb{S}$ für einen Schaltvorgang gemäß \ref{sec:Gangwechsel} interpretiert.

\subsection{Sensitivität des Zustandes $\pmb{x}$ von $k_\mathrm{ss}$}

kann die Sensitivitätsfunktion wie fo 


\section{Fisher-Informationsmatrix}





\chapter{Schätzverfahren}\label{ch:ch3}


Des weiteren kann durch die Betrachtung der Widerstandskräfte $F_\mathrm{r}(v)$ und $F_\mathrm{ad}(v)$ als Eingangsgrößen, das System weiterhin in lineare Form
\begin{equation}\label{eq:sys_linex}
\dot{\pmb{x}}_\mathrm{ex} = \pmb{A}_\mathrm{ex}\,\pmb{x}_\mathrm{ex} + \pmb{B}_\mathrm{ex}\,\pmb{u}_\mathrm{ex}
\end{equation}
angegeben werden. Hierfür werden die beiden Widerstandskräfte in \eqref{eq:dynwc} durch $T^*_\mathrm{res}$ ersetzt, sodass sich 
\begin{equation}\label{eq:dynwc_lin}
\dot{\omega}_\mathrm{C} = \left[ 2\,\left(k_\mathrm{ss}\left(\frac{\phi_2}{i_{\mathrm{D}}} - \phi_\mathrm{C}\right) + d_\mathrm{ss}\left(\frac{\omega_2}{i_{\mathrm{D}}} - \omega_\mathrm{C}\right)\right) - T^*_\mathrm{res} - r_\mathrm{dyn}\, F_\mathrm{g}\right]/I^\mathrm{eff}_\mathrm{C}.
\end{equation}
 ergibt. Der Eingangsvektor des linearen Systems ist dann definiert als 
\begin{equation}
\pmb{u}_\mathrm{ex} = \begin{bmatrix} \mathrm{T}_{\mathrm{In}} & T^*_\mathrm{res} & \mathrm{T}_{\mathrm{K81}} & \mathrm{T}_{\mathrm{K38}} & \mathrm{T}_{\mathrm{B08}} & \mathrm{T}_{\mathrm{B05}} & \mathrm{T}_{\mathrm{B06}} & \mathrm{T}_{\mathrm{K27}} & \zeta \end{bmatrix}^T.
\end{equation}

Die Funktion $\pmb{f}(\pmb{x}_\mathrm{ex})$ und die Eingangmatrix $\pmb{B}_\mathrm{ex,nl}$ des nichtlinearen Systems, sowie die Systemmatrix $\pmb{A}_\mathrm{ex}$ und die Eingangmatrix $\pmb{B}_\mathrm{ex}$ des linearen Systems sind im Anhang mit numerischen Werten für die Zähnezahlen der Zahnräder und der Massenträgheitsmomente angegeben.



Mit den eingesetzten numerischen Werten der Zahnräder, der Massenträgheitsmomente, der Übersetzung der Differentials und dem dynamischen Radius der Räder ergeben sich die Systemmatrizen zu
\begin{align}
\pmb{A}_\mathrm{ex,23} &= \begin{bmatrix} 0 & -0,799\,d_\mathrm{ss} & -1,97\,k_\mathrm{ss} & 1,97\,d_\mathrm{ss} \\
0 & -0,248\,d_\mathrm{ss} & -0,612\,k_\mathrm{ss} & 0,612\,d_\mathrm{ss}\\
0 & 0,405 & 0 & -1,0\\
0 & \frac{0,81\,d_\mathrm{ss}}{m_\mathrm{veh}\,0,108 + 6,4} & \frac{2,0\,k_\mathrm{ss}}{m_\mathrm{veh}\,0,108 + 6,4} & \frac{-2,0\,d_\mathrm{ss}}{m_\mathrm{veh}\,0,108 + 6,4})\end{bmatrix}\\
\nonumber \\
\pmb{B}_\mathrm{ex,23} &= \begin{bmatrix} 10,7 & -2,79 & -17,0 & 0 & 0\\
2,44 & 0,0135 & -2,41 & 0 & 0 \\
0 & 0 & 0 & 0 & 0\\
0 & 0 & 0 & \frac{-1}{0,108\,m_\mathrm{veh} + 6,4} & \frac{-3,22\,m_\mathrm{veh}}{0,108\,m_\mathrm{veh} + 6,4} \end{bmatrix}
\end{align}

\section{Kalman-Filter}
\section{Extented Kalman-Filter}
\section{Unknown Input Observer}
\section{Extended Unknown Input Observers}





\chapter{Zusammenfassung und Ausblick}\label{ch:ch4}
Zusammenfassung und Ausblick werden hier beschrieben.
\cleardoublepage

\appendix
\chapter{Anhang}\label{ch:appCh1}
Im Anhang werden Ergebnisse aufgeführt, die zwar im Kontext der Arbeit wesentlich, jedoch nicht für das Verständnis der Arbeit notwendig sind. Beispielsweise sind dies detailliertere Messergebnisse, Herleitungen oder Datenblätter. 

% List of tables and list of figures at end
\listoffigures
\listoftables

% Literaturverzeichnis
% classical bibtex
%\bibliography{references}

% modern biblatex
\printbibliography[heading=bibintoc]

% in any case ...
\label{ch:bib}



% Erklaerung des Autors, netig fuer Bachelor-/Master-Arbeiten der Universitaet Stuttgart
% Declaration of authorship, necessary for bachelor and master thesis at University of Stuttgart
\DeclarationOfAuthorship
\end{document}
